\chapter{绪论}

\section{研究背景与意义}

管道输送系统广泛应用于能源、化工与市政基础设施等场景，其运行安全直接关系到生产连续性、环境安全与维护成本。实际运行过程中，泄漏故障具有隐蔽性强、传播范围广、发现滞后等特点，一旦未被及时发现，容易造成介质损失、设备损坏甚至安全事故。因此，构建能够面向原始传感信号进行自动识别与定位的泄漏检测方法，具有较强的工程应用价值。

传统泄漏检测方法通常依赖人工经验、阈值规则或频域特征分析。此类方法在结构简单、噪声可控的场景下具有一定效果，但面对复杂工况、长时序原始信号以及不同泄漏形态时，往往存在特征设计依赖人工、泛化能力有限、定位精度不足等问题。随着深度学习在时序建模中的应用逐渐成熟，直接从原始信号中学习判别特征，为管道泄漏检测与定位提供了新的技术路径。

基于此，本文围绕双通道管道传感信号，设计了一种两阶段泄漏检测框架。该框架首先对分段时序信号进行泄漏形态分类，再对泄漏样本进行距离估计，从而将复杂定位问题拆分为“状态识别”与“位置估计”两个相对清晰的子任务，以提升模型训练稳定性与整体任务可解释性。

\section{研究问题与本文工作}

结合现有数据与代码实现，本文关注的核心问题可以概括为：如何利用双通道长时序传感信号，在较少人工特征设计的前提下，实现管道泄漏状态识别与泄漏距离估计。

围绕上述问题，本文完成了以下工作：

\begin{enumerate}
  \item 设计了原始 CSV 信号到训练样本的自动化构建流程。针对原始双通道数据，按照固定时间窗口完成切分，并分别构建分类任务与回归任务所需的数据清单。
  \item 构建了两阶段学习任务。其中，第二阶段任务用于识别信号所属类别，第一阶段任务用于对泄漏距离进行回归估计。
  \item 设计了适用于长时序一维信号的卷积神经网络模型。模型采用多层一维卷积模块逐步降采样，并通过全局池化与任务头实现分类或回归输出。
  \item 搭建了统一训练框架，并在分类任务中补充片段级与原始文件级评测指标，包括准确率、macro-F1 和混淆矩阵；在回归任务中除平均绝对误差与均方根误差外，进一步补充误差方向、容差命中率、按距离档位误差和文件级聚合误差分析，为实验结果解释提供更加完整的依据。
\end{enumerate}

\section{论文结构安排}

全文共分为五章，各章内容安排如下：

\begin{enumerate}
  \item 第一章为绪论，介绍研究背景、研究意义、研究问题与本文主要工作。
  \item 第二章介绍管道泄漏探测的基本原理、泄漏声信号形成机制、信号数据特征以及数据集构建过程。
  \item 第三章介绍本文所用深度学习方法基础、两阶段泄漏检测模型的整体结构、网络设计与训练方法。
  \item 第四章介绍实验设计、评价指标与结果分析，重点讨论分类性能、距离估计误差、统计基线对比及影响因素。
  \item 第五章总结全文工作，并对后续改进方向进行展望。
\end{enumerate}

\chapter{管道泄漏探测原理与数据集构建}

\section{管道泄漏探测基本原理}

管道泄漏检测的核心目标是从传感器采集到的压力、振动或声学响应中识别异常状态，并进一步推断泄漏位置。管道处于正常运行状态时，内部流体流动与外部结构响应通常保持相对稳定；当管壁出现孔洞、裂纹或连接处密封失效时，流体在压差作用下从破损位置喷出，局部流场会发生突变，并激发出与泄漏有关的瞬态波动和持续扰动。这些扰动沿管壁和流体介质传播，被布置在不同位置的传感器接收后表现为可观测的时序信号。

从信号分析角度看，泄漏探测可以理解为从原始时序信号中提取能够区分正常状态、不同泄漏形态以及泄漏位置的特征。传统方法通常依赖时域统计量、频域能量分布或人工设计的阈值规则；而深度学习方法则尝试直接从原始波形中学习特征表示。本文采用后一种思路，以固定长度的分段信号为输入，通过一维卷积网络自动提取局部波形模式和跨通道差异，从而完成泄漏识别与距离估计。

\section{泄漏声信号形成机制}

泄漏声信号的形成主要来源于泄漏点附近流体状态的突变。当管道内部介质经泄漏孔口向外喷射时，局部区域会产生湍流、压力脉动和结构振动。这些扰动一方面沿流体传播，另一方面也会耦合到管壁结构中形成振动响应。对于布置在管道不同位置的左右两个传感通道而言，同一泄漏事件到达两个传感位置的传播路径、衰减程度和相位关系并不完全相同，因此双通道信号中包含了与泄漏位置相关的信息。

泄漏距离估计的物理依据正是这种传播差异。若泄漏点靠近某一侧传感器，该侧通道可能表现出更明显的幅值变化、更早的响应或不同的局部频率成分；若泄漏点位于两侧传感器之间的不同位置，左右通道之间的相关性、能量差异和波形形态也会随之变化。本文没有手工构造到达时间差或频域峰值等特征，而是保留左右通道原始分段信号，由卷积网络在训练过程中学习这些隐含差异。

\section{数据特征与建模需求}

本文数据具有以下几个特点。第一，单段信号长度较长。采样率为 25600 Hz、窗口长度为 5 s 时，每个通道包含 128000 个采样点，直接使用全连接网络建模会带来较高参数量和过拟合风险。第二，信号为典型一维时序数据，局部邻域内的波形变化具有连续性，因此适合采用一维卷积提取局部模式。第三，回归任务使用左右双通道输入，通道间差异本身具有定位意义，模型需要同时利用单通道局部特征和双通道联合特征。第四，距离标签虽然以回归形式使用，但实际取值为若干离散档位，因此实验分析不能只看总体平均误差，还应关注不同距离档位上的误差差异。

基于上述特征，本文在数据处理时采用固定窗口切分和逐段标准化，在模型设计时采用多层一维卷积逐步降采样，在结果分析时同时报告总体指标、按原始文件聚合后的指标以及按距离档位拆分的误差表现。

\section{原始数据组织方式}

本项目原始数据以 CSV 文件形式存储，按照是否泄漏划分为两个目录，即 \texttt{raw/leak} 与 \texttt{raw/normal}。每个文件对应一组实验采样记录，文件名采用 \texttt{ABC} 编号形式，例如 \texttt{ABC91.csv}。从代码实现上看，原始文件被统一读入，并在跳过首行后转为数值矩阵。数据样本中包含左右两个传感通道，分别对应后续建模中使用的左通道信号与右通道信号。

这种数据组织方式有两个优点。其一，原始信号与实验编号一一对应，便于后续依据编号映射标签；其二，双通道结构保留了泄漏传播在不同传感位置上的响应差异，为后续距离估计提供了基础。

\section{标签映射与任务拆分}

数据准备脚本 \texttt{src/leak\_detection/cli/prepare\_dataset.py} 中定义了编号到标签的映射规则。对于发生泄漏的实验样本，每个编号区间同时对应两个标签信息：一是泄漏距离 \texttt{distance}，二是泄漏形态类别 \texttt{shape}；对于正常样本，则只赋予正常类别标签，不参与距离回归任务。

基于这一映射关系，本文将原始问题拆分为两个子任务：

\begin{enumerate}
  \item \texttt{Stage2} 为分类任务。输入为单通道分段信号，输出为三类标签，用于区分正常状态与两类不同泄漏形态。
  \item \texttt{Stage1} 为回归任务。输入为同一时间窗口下的左右双通道信号对，输出为泄漏距离，用于估计泄漏位置。
\end{enumerate}

这种两阶段设计的考虑在于：如果直接从原始信号同时完成状态识别与位置估计，模型需要同时学习类别边界与位置回归关系，训练难度较大。先完成状态判别，再针对泄漏样本估计距离，可以在一定程度上降低任务耦合度，使模型结构与实验分析更加清晰。

\section{分段策略与样本生成}

在数据处理过程中，原始长时序信号按照固定窗口进行切分。项目中默认采样率为 25600 Hz，分段时长为 5 s，因此单段样本长度为：

\begin{equation}
L = 25600 \times 5 = 128000
\end{equation}

对于每个原始文件，脚本分别读取左通道与右通道数据，并按不重叠窗口进行切段。每个时间窗口都被保存为独立 CSV 文件，从而得到可直接供模型读取的分段样本。对于分类任务，左右通道被视为两条独立单通道样本；对于回归任务，左右通道在相同时间窗口下配对，形成双通道输入。

这种固定窗口切分方式具有实现简单、样本规模可扩展、便于批量训练等优点。同时，固定长度输入也便于卷积网络进行统一建模。不过，该策略也存在局限，即不同窗口之间没有显式时序关联，若要进一步利用跨窗口上下文信息，后续可以考虑引入重叠切窗或序列级建模方法。

\section{数据集统计}

根据当前仓库中已经生成的数据清单，本文所使用的数据规模如表\ref{tab:dataset_stats}所示。

\begin{table}[!htbp]
  \centering
  \caption{当前数据集统计结果}
  \label{tab:dataset_stats}
  \vspace{0.5em}
  \begin{tabular}{lcc}
    \toprule
    任务 & 样本数量 & 监督目标 \\
    \midrule
    Stage2 分类 & 3024 & 三分类标签 \\
    Stage1 回归 & 1122 & 泄漏距离 \\
    \bottomrule
  \end{tabular}
\end{table}

进一步统计发现，\texttt{Stage2} 任务的三类样本分布约为正常类 780 条、泄漏类别 1 为 1440 条、泄漏类别 2 为 804 条，存在一定程度的类别不均衡；\texttt{Stage1} 任务的距离标签共有 9 个离散取值，主要集中在 12、14 和 16 等位置附近。上述数据分布特征意味着：分类任务需要关注类别偏置问题，回归任务则需要重点分析不同距离区间上的预测误差。

\chapter{两阶段泄漏检测模型设计}

\section{深度学习方法基础}

本文使用的模型属于监督学习框架。给定输入信号 $x$ 与对应标签 $y$，模型通过参数化函数 $f_{\theta}(x)$ 建立从信号到类别或距离的映射关系，并通过最小化损失函数不断更新参数 $\theta$。对于分类任务，标签为离散类别，训练目标是提高类别判别准确性；对于距离回归任务，标签为连续数值，训练目标是减小预测距离与真实距离之间的偏差。

一维卷积神经网络适用于处理时序信号。卷积核沿时间轴滑动，在局部窗口内共享参数，从而能够提取局部波形变化、短时冲击响应和局部振荡模式。与直接连接所有采样点的全连接网络相比，一维卷积具有参数共享和局部感受野的特点，能够在较少参数量下处理较长序列。随着网络层数增加，后续特征能够覆盖更大的时间范围，从而形成由局部到全局的多层次表示。

在本文任务中，批归一化用于缓解不同 batch 特征分布波动带来的训练不稳定问题，GELU 激活函数用于引入非线性表达能力，自适应平均池化用于将不同时间位置上的高维特征压缩为固定长度的全局表示。优化器采用 AdamW，其在 Adam 自适应学习率更新基础上引入解耦权重衰减，有助于提升训练稳定性并降低过拟合风险。上述模块共同构成了本文模型训练的深度学习基础。

\section{总体框架}

本文采用的模型框架如图示流程所示，可概括为“数据切分、分类识别、距离估计”三个环节。首先，将原始双通道采样信号切分为固定长度片段；其次，利用单通道分类模型对分段信号进行状态识别；最后，对已知为泄漏的信号片段，利用双通道回归模型估计泄漏距离。

从代码实现角度看，任务入口统一由 \texttt{src/leak\_detection/cli/train.py} 调用训练器，训练器根据配置文件自动选择任务类型、数据集类和网络结构。这样的设计使分类任务与回归任务共享训练流程，只在数据组织、损失函数和评价指标上进行区分，有利于后续对比实验与复现实验。

\section{输入表示与归一化}

模型输入由 \texttt{src/leak\_detection/data/segmented.py} 完成读取。对于分类任务，每个样本是形状为 $1 \times 128000$ 的单通道张量；对于回归任务，每个样本由左右通道堆叠形成形状为 $2 \times 128000$ 的双通道张量。

考虑到不同实验样本的信号幅值范围可能存在差异，代码中对每个样本执行逐段标准化处理。设原始信号为 $x$，均值为 $\mu$，标准差为 $\sigma$，则归一化后的信号为：

\begin{equation}
\hat{x} = \frac{x - \mu}{\sigma + 10^{-6}}
\end{equation}

当标准差接近零时，仅进行去均值操作，以避免除零问题。该策略有助于降低量纲差异对训练过程的影响，使网络更关注波形形态与通道间相对变化。

\section{一维卷积特征提取网络}

本文采用的一维卷积主干网络定义于 \texttt{src/leak\_detection/models/signal\_models.py}。网络由若干级卷积模块堆叠而成，每一级模块均包含以下操作：

\begin{enumerate}
  \item 一层带步长的一维卷积，用于完成特征提取与时域降采样；
  \item 批归一化层，用于稳定训练过程；
  \item GELU 激活函数，用于增强非线性表达能力；
  \item 第二层一维卷积与批归一化、GELU 组合，用于进一步提取局部模式。
\end{enumerate}

模型配置中卷积通道数依次设置为 32、64、128 和 256，对应卷积核大小分别为 15、9、7 和 5，步长均为 4。通过多层降采样，长达 128000 点的原始时序被逐步压缩到较短的特征表示，再经过自适应平均池化聚合为全局特征向量。最后，采用全连接层构建任务头。

这一结构的优点在于实现简单、参数量可控，并且适合处理长度较长的一维信号。相比需要更大算力与更复杂位置编码机制的序列模型，一维卷积网络在当前数据规模和实验条件下更易训练，也更适合本科毕业设计阶段的工程实现。

\section{分类头与回归头设计}

\subsection{Stage2 分类模型}

\texttt{Stage2} 任务使用单通道输入，模型主干输出全局特征后，经由两层全连接网络映射到三维类别空间，最终输出对应三类标签的预测结果。分类任务损失函数采用交叉熵损失，其表达式为：

\begin{equation}
\mathcal{L}_{cls} = - \sum_{c=1}^{C} y_c \log \hat{y}_c
\end{equation}

其中，$C$ 表示类别数量，$y_c$ 为真实标签的 one-hot 编码，$\hat{y}_c$ 为模型输出的类别概率。

\subsection{Stage1 回归模型}

\texttt{Stage1} 任务使用双通道输入，模型输出单一实数作为泄漏距离预测值。为增强对异常误差的鲁棒性，当前实现中优先采用 Smooth L1 损失：

\begin{equation}
\mathcal{L}_{reg}(d,\hat{d}) =
\left\{
\begin{array}{ll}
\frac{1}{2}(d-\hat{d})^2, & |d-\hat{d}| < 1 \\
|d-\hat{d}| - \frac{1}{2}, & |d-\hat{d}| \geq 1
\end{array}
\right.
\end{equation}

其中，$d$ 为真实距离，$\hat{d}$ 为预测距离。相较于均方误差，Smooth L1 在误差较大时对离群点的敏感性更低，有利于提升回归训练的稳定性。

\section{训练策略}

统一训练流程定义于 \texttt{src/leak\_detection/training/trainer.py}。当前实验采用 AdamW 优化器，初始学习率为 0.001，权重衰减为 0.0001，并使用余弦退火学习率调度器。训练最大轮数设置为 40，早停轮数设置为 8，当验证集指标长期未提升时提前终止训练。

此外，为避免梯度爆炸对长时序训练造成不稳定影响，训练过程中对梯度范数进行裁剪，阈值为 1.0。对分类任务，验证指标选择准确率；对回归任务，模型选择仍以验证集 MAE 为主，同时在测试阶段补充 RMSE、误差方向、容差命中率、最大误差、最小误差以及文件级聚合指标。训练完成后，程序自动加载最佳验证模型，并在测试集上输出最终指标。

\chapter{实验设计与结果分析}

\section{实验环境与配置}

本文实验基于 Python 与 PyTorch 框架完成，数据集构建与训练流程均通过命令行工具调用。当前项目提供了两个主要配置文件：\texttt{configs/stage2.yaml} 用于分类实验，\texttt{configs/stage1.yaml} 用于回归实验。两者在主干网络结构上保持一致，仅在输入形式、任务头与评价指标上有所区别。

从实验复现角度出发，本文建议固定随机种子为 42，保持训练集、验证集和测试集划分一致，以保证不同实验之间具有可比性。若后续进行超参数对比实验，可分别围绕卷积通道数、卷积核大小、损失函数和窗口长度展开。

\section{评价指标}

\subsection{分类任务指标}

对于 \texttt{Stage2} 分类任务，本文采用准确率、macro-F1 和混淆矩阵作为主要评价方式。准确率定义为：

\begin{equation}
Accuracy = \frac{N_{correct}}{N_{total}}
\end{equation}

其中，$N_{correct}$ 表示预测正确的样本数量，$N_{total}$ 表示测试样本总量。考虑到分类任务存在一定类别不均衡，本文进一步引入 macro-F1 指标。该指标先分别计算每个类别的 F1 值，再对各类别结果进行等权平均，因此能够更好地反映少数类识别效果。此外，本文同时统计片段级混淆矩阵与原始文件级混淆矩阵，用于分析不同类别之间的混淆关系。

\subsection{回归任务指标}

对于 \texttt{Stage1} 回归任务，本文采用平均绝对误差与均方根误差作为主要指标，分别表示为：

\begin{equation}
MAE = \frac{1}{N}\sum_{i=1}^{N}|d_i-\hat{d}_i|
\end{equation}

\begin{equation}
RMSE = \sqrt{\frac{1}{N}\sum_{i=1}^{N}(d_i-\hat{d}_i)^2}
\end{equation}

其中，$d_i$ 为真实距离，$\hat{d}_i$ 为预测距离。MAE 反映平均偏差水平，RMSE 对较大误差更加敏感，二者结合能够较全面地反映定位性能。

为增强回归结果的物理可解释性，本文进一步统计有符号误差 $e_i=\hat{d}_i-d_i$、最大绝对误差、最小绝对误差以及容差命中率。容差命中率定义为：

\begin{equation}
Acc_{\tau} = \frac{1}{N}\sum_{i=1}^{N}\mathbf{1}(|\hat{d}_i-d_i|\leq \tau)
\end{equation}

其中，$\tau$ 表示允许误差阈值。本文重点报告 $\tau=1.0$ 时的结果，用于衡量预测值是否落在真实距离附近一个单位范围内。由于同一原始文件会被切分为多个片段，本文还对同一原始文件下的片段预测结果求平均，得到文件级距离预测，并计算对应的文件级 MAE、RMSE 和最大误差。

\section{实验结果}

基于当前仓库中的正式训练结果，本文分别对 \texttt{Stage2} 分类任务与 \texttt{Stage1} 回归任务进行分析。所有实验均采用相同的随机种子与文件级划分方式，以保证训练集、验证集与测试集之间具有可比性。

\subsection{Stage2 分类结果}

\texttt{Stage2} 任务以单通道分段信号为输入，输出正常状态与两类泄漏形态的三分类结果。实验结果如表\ref{tab:stage2_results} 所示。可以看出，模型在测试集上取得了极高的片段级分类性能，测试集片段级准确率达到 0.9968，片段级 macro-F1 达到 0.9972。进一步将同一原始文件下的所有片段预测进行聚合后，文件级准确率与文件级 macro-F1 均达到 1.0000，说明模型在当前随机文件级划分下，对未见测试文件具备很强的识别能力。

\begin{table}[!htbp]
  \centering
  \caption{Stage2 分类任务测试结果}
  \label{tab:stage2_results}
  \vspace{0.5em}
  \begin{tabular}{lcc}
    \toprule
    评测层级 & 指标 & 数值 \\
    \midrule
    片段级 & Accuracy & 0.9968 \\
    片段级 & Macro-F1 & 0.9972 \\
    文件级 & Accuracy & 1.0000 \\
    文件级 & Macro-F1 & 1.0000 \\
    \bottomrule
  \end{tabular}
\end{table}

从混淆矩阵看，片段级测试样本中仅存在一次误分类，即真实类别 1 的一个片段被误判为类别 2，其余样本均被正确识别。文件级混淆矩阵则全部位于主对角线上，表明测试集中 26 个原始文件均被正确分类。上述结果不仅说明总体准确率较高，也说明该结果并非仅由多数类样本支撑，而是在各类别上都表现稳定。

\subsection{Stage1 距离回归结果}

\texttt{Stage1} 任务以同一时间窗口下的左右双通道信号对为输入，输出泄漏距离预测值。实验结果如表\ref{tab:stage1_results} 所示。模型在验证集上的最佳平均绝对误差为 0.3148；在测试集上，平均绝对误差为 0.3703，均方根误差为 0.5934。测试误差相较验证集略有上升，但整体差距较小，表明模型在距离估计任务上具有较好的稳定性。

\begin{table}[!htbp]
  \centering
  \caption{Stage1 回归任务结果}
  \label{tab:stage1_results}
  \vspace{0.5em}
  \begin{tabular}{lcc}
    \toprule
    数据划分 & 指标 & 数值 \\
    \midrule
    验证集最优模型 & MAE & 0.3148 \\
    测试集 & MAE & 0.3703 \\
    测试集 & RMSE & 0.5934 \\
    \bottomrule
  \end{tabular}
\end{table}

若与简单统计基线进行比较，本文模型表现出明显优势。本文选取两种不使用输入信号的统计基线：一是始终预测训练集距离均值，二是始终预测训练集距离中位数。对比结果如表\ref{tab:stage1_baseline} 所示。可以看出，统计基线的测试集 MAE 分别为 4.4040 和 3.8571，显著高于本文模型的 0.3703。这说明模型并非仅仅学习到了距离标签的整体分布，而是有效捕获了双通道信号模式与泄漏距离之间的映射关系。

\begin{table}[!htbp]
  \centering
  \caption{Stage1 回归任务统计基线对比}
  \label{tab:stage1_baseline}
  \vspace{0.5em}
  \begin{tabular}{lccc}
    \toprule
    方法 & 预测策略 & MAE & RMSE \\
    \midrule
    训练集均值基线 & 恒定预测 9.3594 & 4.4040 & 4.7855 \\
    训练集中位数基线 & 恒定预测 11.0000 & 3.8571 & 4.9425 \\
    本文一维卷积模型 & 双通道信号回归 & 0.3703 & 0.5934 \\
    \bottomrule
  \end{tabular}
\end{table}

进一步从物理误差角度分析，本文统计了最小绝对误差、最大绝对误差、误差方向和预测距离范围，结果如表\ref{tab:stage1_physical_error} 所示。片段级最小绝对误差为 0.0060，最大绝对误差为 2.4668；有符号误差最小值为 -2.4668，最大值为 1.1769，说明较大误差主要表现为低估。片段级误差不超过 1 个距离单位的比例为 0.9286。将同一原始文件下多个片段的预测取均值后，文件级 MAE 降至 0.2534，RMSE 降至 0.4603，说明多片段聚合能够进一步削弱单段预测波动。

\begin{table}[!htbp]
  \centering
  \caption{Stage1 回归任务物理误差分析}
  \label{tab:stage1_physical_error}
  \vspace{0.5em}
  \begin{tabular}{lcc}
    \toprule
    指标 & 片段级 & 文件级 \\
    \midrule
    MAE & 0.3703 & 0.2534 \\
    RMSE & 0.5934 & 0.4603 \\
    最小绝对误差 & 0.0060 & 0.0029 \\
    最大绝对误差 & 2.4668 & 1.8852 \\
    平均有符号误差 & -0.1504 & -0.1504 \\
    最大低估误差 & -2.4668 & -1.8852 \\
    最大高估误差 & 1.1769 & 0.3124 \\
    误差 $\leq 1.0$ 占比 & 0.9286 & 0.9524 \\
    预测距离范围 & 1.9274--16.1776 & 2.0277--15.7598 \\
    \bottomrule
  \end{tabular}
\end{table}

由于距离标签为多个离散档位，本文进一步按真实距离拆分误差，结果如表\ref{tab:stage1_distance_error} 所示。可以看出，大多数距离档位的 MAE 均低于 0.5，且误差不超过 1 个距离单位的比例接近或达到 1.0；但距离 16 的 MAE 为 1.1186，平均有符号误差为 -1.0627，明显存在低估现象。该现象说明模型在高距离档位上的定位能力相对较弱，可能与该距离档位样本数量、信号衰减特征或与相邻距离档位的波形相似性有关。

\begin{table}[!htbp]
  \centering
  \caption{Stage1 按距离档位的片段级误差}
  \label{tab:stage1_distance_error}
  \vspace{0.5em}
  \begin{tabular}{ccccc}
    \toprule
    真实距离 & 样本数 & MAE & 最大绝对误差 & 平均有符号误差 \\
    \midrule
    2 & 6 & 0.0750 & 0.1482 & 0.0277 \\
    3 & 24 & 0.2519 & 0.8743 & 0.0756 \\
    4 & 6 & 0.2327 & 0.3438 & -0.0435 \\
    5 & 6 & 0.1849 & 0.4042 & -0.1465 \\
    11 & 18 & 0.4296 & 1.3331 & 0.0747 \\
    12 & 30 & 0.3519 & 2.4668 & -0.1672 \\
    14 & 24 & 0.2476 & 0.6062 & -0.1403 \\
    16 & 12 & 1.1186 & 2.2493 & -1.0627 \\
    \bottomrule
  \end{tabular}
\end{table}

从最大误差样本看，片段级最大绝对误差来自 \texttt{ABC155}，真实距离为 12，预测值为 9.5332，误差为 -2.4668；文件级最大误差来自 \texttt{ABC162}，真实距离为 16，多个片段平均预测为 14.1148，误差为 -1.8852。这进一步说明后续优化应重点关注高距离档位及个别异常文件的误差来源。

\begin{table}[!htbp]
  \centering
  \caption{两阶段任务主要实验结果汇总}
  \label{tab:main_results}
  \vspace{0.5em}
  \begin{tabular}{lccc}
    \toprule
    任务 & 测试指标 1 & 测试指标 2 & 备注 \\
    \midrule
    Stage2 分类 & Accuracy = 0.9968 & Macro-F1 = 0.9972 & 文件级 Accuracy = 1.0000 \\
    Stage1 回归 & MAE = 0.3703 & RMSE = 0.5934 & 双通道距离估计 \\
    \bottomrule
  \end{tabular}
\end{table}

\section{结果分析}

综合分类与回归实验结果可以看出，本文提出的两阶段框架在当前数据集与划分设置下取得了较好的性能。对于 \texttt{Stage2} 分类任务，模型几乎实现了完美识别，说明基于一维卷积的特征提取方式能够较好地捕获分段信号中的泄漏形态差异。特别是在引入原始文件级聚合评测后，模型依然获得 1.0000 的文件级准确率与文件级 macro-F1，这表明该结果并非单纯由同一文件下多个相似片段造成的片段级虚高。

对于 \texttt{Stage1} 回归任务，测试集 MAE 为 0.3703，RMSE 为 0.5934，明显优于简单均值预测或中位数预测等统计基线，表明双通道输入确实为泄漏距离估计提供了有效信息。由于当前距离标签为多个离散取值，该误差水平说明模型通常能够将预测结果控制在相邻距离档位附近，具有较好的距离估计能力。

从模型性能分析角度看，本文当前实验已经完成了三个层面的对比。第一，与不使用信号输入的统计基线相比，一维卷积模型显著降低了 MAE 和 RMSE，验证了信号特征对于距离估计的必要性。第二，片段级与文件级结果对比表明，同一原始文件的多片段预测均值能够进一步降低误差，说明在实际应用中可优先采用文件级或时间段级聚合结果作为最终定位输出。第三，按距离档位拆分后的误差表明，总体平均指标不能完全反映模型行为，高距离档位 16 的低估问题是当前回归任务的主要误差来源。

消融实验方面，本文当前代码已保留若干可配置项，例如回归损失可在 Smooth L1 与 MSE 之间切换，网络通道数、卷积核大小、窗口长度和归一化策略也可以作为后续变量。受当前实验周期限制，本文未将所有消融结果作为正式表格展开，而是将已完成的统计基线和聚合层级对比作为现阶段模型有效性验证。后续若进一步完善实验，可围绕“是否使用双通道输入”“Smooth L1 与 MSE 损失对比”“不同窗口长度对定位精度的影响”“是否进行文件级聚合”四个方向开展系统消融，以更细致地量化各模块贡献。

需要指出的是，本文结果建立在当前随机文件级划分的基础上。虽然训练集、验证集与测试集之间不存在原始文件重叠，但测试集规模仍然相对有限，且不同文件之间可能仍具有较强的工况相似性。因此，本文将当前结果视为现阶段实验条件下的有效结论。后续若需进一步增强说服力，可在保持现有评测指标的基础上，设计更加严格的数据切分策略，并补充跨工况或跨批次泛化实验。

总体而言，当前实验结果验证了两阶段方案的有效性：分类任务能够可靠地区分正常状态与泄漏形态，回归任务能够在较小误差范围内估计泄漏距离。这为后续继续开展模型改进、切分优化和部署适配奠定了较为扎实的实验基础。

\chapter{结论与展望}

\section{全文总结}

本文围绕双通道管道传感信号的泄漏检测问题，完成了数据构建、任务拆分、模型设计与训练流程搭建。针对原始 CSV 时序数据，本文采用固定窗口切分方式构建了分类任务与回归任务所需的数据集，并提出了一个两阶段泄漏检测框架。该框架使用一维卷积神经网络对长时序信号进行建模，其中分类任务用于识别正常状态与泄漏形态，回归任务用于估计泄漏距离。

从工程实现角度看，本文工作已形成较为完整的实验流水线，包括原始数据自动切分、数据加载、模型训练、验证集选择与测试集评估等模块。该流程结构清晰、复现方便，为后续进一步开展对比实验与性能优化提供了基础。

\section{后续工作展望}

尽管本文已经完成了两阶段泄漏检测方法的基本实现，但仍有若干方向值得进一步研究：

\begin{enumerate}
  \item 在数据层面，可进一步扩充不同工况、不同噪声条件和更多距离位置的实验样本，以提升模型泛化能力。
  \item 在模型层面，可尝试引入多尺度卷积、残差结构、注意力机制或时频联合建模方法，以增强对复杂泄漏模式的识别能力。
  \item 在任务层面，可考虑将分类与定位设计为联合学习框架，探索共享特征表示是否能够进一步提升整体性能。
  \item 在应用层面，可围绕实时推理、模型轻量化和部署适配展开研究，为后续实际工程部署提供支持。
\end{enumerate}

总体而言，本文完成的工作验证了基于分段时序信号与深度学习方法开展管道泄漏检测与定位的可行性，为后续进一步研究打下了基础。
